% answers_short.tex —— 2015 年·全国Ⅰ卷·文科 · 简短答案速查
% 由 tools/make_short_answers.py 自 answers.tex 的 \gksolhead/\gkanswer 提取，
% 并与 2015/2015_数学真题_全国1文_答案要点.md 逐题比对后生成，勿手改（改源 MD 后重跑）。

\begin{gknotice}
\textbf{使用说明}：本节只给最终答案便于快速核对，完整推导见解析版。答案为编写者独立推导，\textbf{非官方标准答案}。
\end{gknotice}

\gkpart{选择题}
\begin{center}\begin{tabular}{ccccccccccccc}
\toprule
题号 & 1 & 2 & 3 & 4 & 5 & 6 & 7 & 8 & 9 & 10 & 11 & 12 \\
\midrule
答案 & D & A & C & C & B & B & B & D & C & A & B & C \\
\bottomrule
\end{tabular}\end{center}

\gkpart{填空题}
\begin{itemize}[leftmargin=2.6em]
  \item[\textbf{13.}] $6$
  \item[\textbf{14.}] $1$
  \item[\textbf{15.}] $4$
  \item[\textbf{16.}] $12\sqrt6$
\end{itemize}

\gkpart{解答题}
\begin{itemize}[leftmargin=2.6em]
  \item[\textbf{17.}] (1) $\frac14$；(2) $1$
  \item[\textbf{18.}] (1) 两平面垂直，见证明；(2) $3+2\sqrt5$
  \item[\textbf{19.}] (1) $y=c+d\sqrt x$；(2) $\widehat y=100.6+68\sqrt x$；(3) ① $576.6$ t、$66.32$ 千元；② $x=46.24$ 千元
  \item[\textbf{20.}] (1) $\frac{4-\sqrt7}{3}<k<\frac{4+\sqrt7}{3}$；(2) $2$
  \item[\textbf{21.}] (1) $a\le0$ 时无零点，$a>0$ 时恰有一个零点；(2) 不等式成立，见证明
  \item[\textbf{22.}] (1) $DE$ 为切线，见证明；(2) $60^\circ$
  \item[\textbf{23.}] (1) $C_1:\rho\cos\theta=-2$；$C_2:\rho^2-2\rho\cos\theta-4\rho\sin\theta+4=0$；(2) U：原文未定义三角形顶点 $C_2$；若指圆心则面积为 $\frac12$（条件性结论）
  \item[\textbf{24.}] (1) $(\frac23,2)$；(2) $a>2$
\end{itemize}
